% !PRINTED 121
\weberpage{121}
If now $e$ is an arbitrary element of $Q$, then $Qe = Q$,
and consequently also $Ra_iQe = Ra_iQ$. From this it follows that by
composition with $e$ the totality of the $h_i$ elements of which
$P_i$ is composed is not altered, and that they are therefore
only interchanged among one another by the permutation group $\Pi$.
But one also sees immediately that these $h_i$ elements $B$ are \gesperrt{transitively} connected by $\Pi$; for $P_1$ consists precisely
of all those complexes $B$ which can be derived from one among them by
application of elements of $Q$, and
hence each of these complexes $B$ of $P_i$ can be carried over into
every other by permutations of $\Pi$.

We wish to formulate the result obtained as a theorem
in the following way:

3.\ \gesperrt{If $P$ is a group of degree $n$, $R$ a divisor
of $P$ of index $j$, $Q$ a divisor of $P$ of
degree $q$, then with $Q$ a permutation group $\Pi$
of $j$ figures is simply or multiply isomorphic. The
$j$ figures split into rows of $h_1, h_2, h_3, \ldots$
figures each, transitively connected by $\Pi$, so that}
\weq{11}{j = h_1 + h_2 + h_3 + \cdots}
\gesperrt{and $h_1, h_2, h_3, \ldots$ are divisors of $q$. The
numbers $h_1, h_2, h_3, \ldots$ are the indices of divisors
of $Q$ which result as intersections of $Q$ with
the conjugate groups $a^{-1}Ra$.}

If the various groups $Q_1, Q_2, \ldots$ have, besides the
unit element, no common divisor, which with simple groups $P$ always occurs, the isomorphism is simple.

\begin{center}
\S.\,29.

\gesperrt{The first \gesperrt{Sylow} theorem.}
\end{center}

It is not easy, on the direct path taken by \gesperrt{Cayley},
to go beyond the lowest numbers of degree in the
formation of groups. For further investigations
in this direction, a theorem is of great use, which in a restricted
scope originates from \gesperrt{Cauchy}, but in general
was proved by \gesperrt{Sylow}. This important theorem we shall
%CONT>>
